% Source: https://artofproblemsolving.com/community/c3h2124207_tangent_line_to_circle_proof
% Stable identity: aops-topic-2124207-q01
\textbf{(a)} The common points have $y=mx+b$ and $x^2+(mx+b)^2=r^2$, that is,
\[\left(1+m^2\right)x^2+2mbx+\left(b^2-r^2\right)=0 .\]
Since $1+m^2>0$, this is a genuine quadratic, and the line meets the circle in exactly one point exactly when the quadratic has exactly one real root, i.e. when its discriminant is $0$:
\[4m^2b^2-4\left(1+m^2\right)\left(b^2-r^2\right)=0\iff -4b^2+4r^2\left(1+m^2\right)=0\iff r^2\left(1+m^2\right)=b^2 .\]
In particular $b^2=r^2(1+m^2)>0$, so $b\ne0$.

\textbf{(b)} The double root is
\[x_T=-\dfrac{2mb}{2\left(1+m^2\right)}=-\dfrac{mb}{1+m^2}=-\dfrac{mb\,r^2}{b^2}=-\dfrac{r^2m}{b},\]
using $1+m^2=\dfrac{b^2}{r^2}$ from (a). Then
\[y_T=mx_T+b=\dfrac{-r^2m^2+b^2}{b}=\dfrac{-r^2m^2+r^2+r^2m^2}{b}=\dfrac{r^2}{b}.\]

\begin{asy}
usepackage("amsmath");
size(5.5cm);
defaultpen(fontsize(10pt));
real r=2; pair O=(0,0); pair T=r*dir(60);
real m=-1/sqrt(3); real bb=r^2/T.y;
draw((-3,0)--(5.2,0),Arrow(6)); draw((0,-2.6)--(0,3.6),Arrow(6));
label("$x$",(5.2,0),E); label("$y$",(0,3.6),N);
draw(circle(O,r),linewidth(1));
path tl=(-1.6,m*(-1.6)+bb)--(5,m*5+bb);
draw(tl,blue+linewidth(1));
draw(O--T,red+linewidth(0.9));
pair u=unit(T-O), v=(1,m)/abs((1,m));
draw(T-0.28*u--T-0.28*u+0.28*v--T+0.28*v,red);
dot(O); dot(T); dot((0,bb));

label("$r$",(O+T)/2,W,red); label("$y=mx+b$",(3.2,m*3.2+bb),NE,blue+fontsize(9pt));
label("$O$",O,SW); label("$T$",T,NE); label("$(0,b)$",(0,bb),W,fontsize(9pt));
label("\mbox{}",(min(currentpicture,user=true)+max(currentpicture,user=true))/2,black); // keeps the text after the figure black when LaTeX typesets the labels inline
\end{asy}

\textbf{(c)} The vector $(1,m)$ points along the line, and $\overrightarrow{OT}=\left(-\dfrac{r^2m}{b},\dfrac{r^2}{b}\right)$. Their dot product is
\[-\dfrac{r^2m}{b}+m\cdot\dfrac{r^2}{b}=0,\]
so $OT$ is perpendicular to the tangent line. $\blacksquare$
